\subsection{取值于李代数的微分形式的原函数}

命题 15. 设 \(\mathrm{G}\) 是李群，\(\mathrm{M}\) 是 \(\mathrm{C}^r (r\geqslant 2)\) 类流形，\(\alpha\) 是 \(\mathrm{C}^{r - 1}\) 类的 1 次微分形式，定义在 \(\mathrm{M}\) 上且取值于 \(\mathrm{L}(\mathrm{G})\)，满足 \(d\alpha +[\alpha ]^2 = 0\)。假设 \(\mathrm{M}\) 单连通。对所有 \(x\in \mathrm{M}\) 和 \(s\in \mathrm{G}\)，存在唯一的映射 \(f\)，其属于 \(\mathrm{C}^{r - 1}\) 类，定义在 \(\mathrm{M}\) 上且取值于 \(\mathrm{G}\)，使得 \(f(x) = s\) 且 \(f^{-1}.df = \alpha\)。

\(f\) 的唯一性由 § 3, no. 17, Corollary 2 to Proposition 59 以及 \(\mathbf{M}\) 连通这一事实得出。下面证明 \(f\) 的存在性。存在一个开覆盖 \((\mathrm{U}_i)_{i\in \mathbf{I}}\)，它覆盖 \(\mathbf{M}\)，并且对所有 \(i\in \mathbf{I}\)，存在映射 \(g_i: \mathrm{U}_i\to \mathrm{G}\)，它属于 \(\mathrm{C}^{r - 1}\) 类，并且 \(g_i^{-1}.d g_i = \alpha\) 在 \(\mathrm{U}_i\) 上成立（§ 4, no. 6, Theorem 5）。由 § 3, no. 17, Corollary 2 to Proposition 59，\(g_i g_j^{-1}\) 在 \(\mathrm{U}_i\cap \mathrm{U}_j\) 上局部为常值。令 \(g_i g_j^{-1} = g_{ij}\)。令 \(\mathbf{G}_d\) 为带离散拓扑的群 \(\mathbf{G}\)。映射 \(g_{ij}: \mathrm{U}_i\cap \mathrm{U}_j\to \mathrm{G}_d\) 连续，并且 \(g_{ij}g_{jk} = g_{ik}\) 在 \(\mathrm{U}_i\cap \mathrm{U}_j\cap \mathrm{U}_k\) 上成立。由于 \(\mathbf{M}\) 单连通，存在连续映射 \(\lambda_i: \mathrm{U}_i\to \mathrm{G}_d\)，使得 \(g_i g_j^{-1} = \lambda_i\lambda_j^{-1}\) 在 \(\mathrm{U}_i\cap \mathrm{U}_j\) 上成立。令 \(g\) 为从 \(\mathbf{M}\) 到 \(\mathbf{G}\) 的映射，其在 \(\mathrm{U}_i\) 上的限制为 \(\lambda_i^{-1}g_i\)，对所有 \(i\in \mathbf{I}\) 均如此。这个映射属于 \(\mathrm{C}^{r - 1}\) 类，且 \(g^{-1}d g = \alpha\)。映射 \(f\) 从 \(\mathbf{M}\) 到 \(\mathbf{G}\)，由 \(f = s(g(x))^{-1}g\) 定义，并满足条件 \(f^{-1}.df = \alpha\) 和 \(f(x) = s\)。

